\chapter{Boolean Analysis --- Noise Stability}
%%% generated by the blueprint pipeline from TCSlib.BooleanAnalysis.ThresholdFunctions.NoiseStability
\section{Overview}
This chapter defines uniform noise stability for dimension-indexed function classes and records the influence-to-noise reduction, Peres's bound for linear threshold functions, and Kane's bounds for polynomial threshold functions.

\section{Declarations}
\begin{definition}[Dimension-indexed function classes]
\lean{BooleanAnalysis.ThresholdFunctions.BooleanFunctionClass}
\leanok
\uses{BooleanAnalysis.BooleanFunc}
A class of functions on Boolean cubes is a family $B=(B_n)_{n\in\mathbb N}$, where $B_n$ is a set of real-valued functions on $\{0,1\}^{n}$.
\end{definition}

\begin{definition}[Boolean-valued function classes]
\lean{BooleanAnalysis.ThresholdFunctions.IsBooleanValuedClass}
\leanok
\uses{BooleanAnalysis.BooleanFunc,BooleanAnalysis.ThresholdFunctions.BooleanFunctionClass,BooleanAnalysis.isPmOne}
For a family $B=(B_n)_{n\in\mathbb N}$ of sets of functions $\{0,1\}^{n}\to\mathbb R$, being Boolean-valued means that every $f\in B_n$, in every dimension $n$, satisfies $f(x)\in\{-1,1\}$ for all $x\in\{0,1\}^{n}$.
\end{definition}

\begin{definition}[Polynomial threshold classes]
\lean{BooleanAnalysis.ThresholdFunctions.polynomialThresholdClass}
\leanok
\uses{BooleanAnalysis.BooleanFunc,BooleanAnalysis.ThresholdFunctions.IsPolynomialThreshold,BooleanAnalysis.isPmOne}
For $k,n\in\mathbb N$, the degree-at-most-$k$ polynomial threshold class consists of functions $f:\{0,1\}^{n}\to\{-1,1\}$ for which there is a real multilinear polynomial $p$ of degree at most $k$ such that $f(x)=\operatorname{sgn}_{+}(p((-1)^{x_1},\ldots,(-1)^{x_n}))$ for every $x$, where $\operatorname{sgn}_{+}(t)=1$ for $t\geq0$ and $\operatorname{sgn}_{+}(t)=-1$ for $t<0$.
\end{definition}

\begin{definition}[Linear threshold classes]
\lean{BooleanAnalysis.ThresholdFunctions.linearThresholdClass}
\leanok
\uses{BooleanAnalysis.BooleanFunc,BooleanAnalysis.ThresholdFunctions.IsLinearThreshold,BooleanAnalysis.isPmOne}
For every $n\in\mathbb N$, the linear threshold class consists of functions $f:\{0,1\}^{n}\to\{-1,1\}$ admitting real coefficients $c,a_1,\ldots,a_n$ with
\[
f(x)=\operatorname{sgn}_{+}\left(c+\sum_{i=1}^{n}a_i(-1)^{x_i}\right),
\]
where $\operatorname{sgn}_{+}(t)=1$ for $t\geq0$ and $-1$ for $t<0$.
\end{definition}

\begin{definition}[Uniform noise stability]
\lean{BooleanAnalysis.ThresholdFunctions.IsUniformlyNoiseStable}
\leanok
\uses{BooleanAnalysis.BooleanFunc,BooleanAnalysis.ThresholdFunctions.BooleanFunctionClass,BooleanAnalysis.ThresholdFunctions.IsBooleanValuedClass,BooleanAnalysis.ThresholdFunctions.noiseSensitivity}
For a family $B=(B_n)_{n\in\mathbb N}$ of sets of real-valued functions on $\{0,1\}^{n}$, uniform noise stability means that every member is $\{-1,1\}$-valued and there exists a function $\varepsilon:\mathbb R\to\mathbb R$ with $0\leq\varepsilon(\delta)\leq1$ for $0\leq\delta\leq1/2$ and $\lim_{\delta\to0^{+}}\varepsilon(\delta)=0$, such that
\[
\Pr[f(X)\neq f(Y)]\leq\varepsilon(\delta)
\quad(n\in\mathbb N,\ f\in B_n,\ 0<\delta\leq1/2).
\]
Here $X$ is uniform on $\{0,1\}^{n}$ and $Y$ is obtained by flipping each coordinate of $X$ independently with probability $\delta$; no range condition is imposed on $\varepsilon$ outside $[0,1/2]$.
\end{definition}

\begin{definition}[Negated coordinate identification]
\lean{BooleanAnalysis.ThresholdFunctions.identifiedNegation}
\leanok
\uses{BooleanAnalysis.BoolCube,BooleanAnalysis.BooleanFunc}
For every $n,m\in\mathbb N$, a function $f:\{0,1\}^{n}\to\mathbb R$, a vector $z\in\{0,1\}^{n}$ and a map $\pi:\{1,\ldots,n\}\to\{1,\ldots,m\}$, negating according to $z$ and identifying according to $\pi$ gives the function $g:\{0,1\}^{m}\to\mathbb R$ defined by
\[
g(w)=f\bigl((z_i+w_{\pi(i)}\pmod2)_{i=1}^{n}\bigr).
\]
\end{definition}

\begin{definition}[Closure under negation and identification]
\lean{BooleanAnalysis.ThresholdFunctions.IsClosedUnderInputNegationAndIdentification}
\leanok
\uses{BooleanAnalysis.BoolCube,BooleanAnalysis.BooleanFunc,BooleanAnalysis.ThresholdFunctions.BooleanFunctionClass,BooleanAnalysis.ThresholdFunctions.identifiedNegation}
For a family $B=(B_n)_{n\in\mathbb N}$ of sets of real-valued functions on $\{0,1\}^{n}$, closure under input negation and identification means that, for all $n,m\in\mathbb N$, $f\in B_n$, $z\in\{0,1\}^{n}$ and maps $\pi:\{1,\ldots,n\}\to\{1,\ldots,m\}$, the function
\[
w\longmapsto f\bigl((z_i+w_{\pi(i)}\pmod2)_{i=1}^{n}\bigr)
\]
belongs to $B_m$.
\end{definition}

\begin{definition}[Reciprocal block count]
\lean{BooleanAnalysis.ThresholdFunctions.noiseBlockCount}
\leanok
For a real parameter $\delta$, the block count is the natural number
\[
m(\delta)=\max\{0,\lfloor\delta^{-1}\rfloor\},
\]
where the floor is the integer floor and $0^{-1}$ is defined to be $0$; in particular $m(\delta)=0$ for $\delta\leq0$.
\end{definition}

\begin{lemma}[Nonnegative stability at nonnegative correlation]
\lean{BooleanAnalysis.ThresholdFunctions.noiseStability_nonneg}
\leanok
\uses{BooleanAnalysis.BooleanFunc,BooleanAnalysis.ThresholdFunctions.noiseStability,BooleanAnalysis.stability_formula}
\difficulty{3}
For every $n\in\mathbb N$, $f:\{0,1\}^{n}\to\mathbb R$ and real $\rho\geq0$, the noise stability
\[
\operatorname{Stab}_{\rho}(f)=\sum_{S\subseteq\{1,\ldots,n\}}\rho^{|S|}\widehat f(S)^2
\]
is nonnegative, where $\widehat f(S)=2^{-n}\sum_{x\in\{0,1\}^{n}}f(x)\prod_{i\in S}(-1)^{x_i}$ is the normalized Walsh coefficient (and $\rho^0=1$, including at $\rho=0$).
\end{lemma}

\begin{lemma}[Half bound for totalized sensitivity]
\lean{BooleanAnalysis.ThresholdFunctions.noiseSensitivity_le_half}
\leanok
\uses{BooleanAnalysis.BooleanFunc,BooleanAnalysis.ThresholdFunctions.noiseSensitivity,BooleanAnalysis.ThresholdFunctions.noiseStability,BooleanAnalysis.ThresholdFunctions.noiseStability_nonneg}
\difficulty{3}
For every $n\in\mathbb N$, real-valued $f:\{0,1\}^{n}\to\mathbb R$ and real $\delta\leq1/2$, including negative $\delta$, the total real-valued noise sensitivity satisfies
\[
\operatorname{NS}_{\delta}(f)=\frac{1-\sum_{S\subseteq\{1,\ldots,n\}}(1-2\delta)^{|S|}\widehat f(S)^2}{2}\leq\frac12,
\]
where $\widehat f(S)=2^{-n}\sum_{x\in\{0,1\}^{n}}f(x)\prod_{i\in S}(-1)^{x_i}$.
\end{lemma}

\begin{lemma}[Bounds on the reciprocal block count]
\lean{BooleanAnalysis.ThresholdFunctions.noiseBlockCount_spec}
\leanok
\uses{BooleanAnalysis.ThresholdFunctions.noiseBlockCount}
\difficulty{4}
For every real $\delta$ with $0<\delta\leq1/2$, the natural number $m=\lfloor1/\delta\rfloor$ satisfies
\[
2\leq m,\qquad m\leq\frac1\delta,\qquad \frac1{2\delta}\leq m.
\]
\end{lemma}

\begin{lemma}[Square-root bound per block]
\lean{BooleanAnalysis.ThresholdFunctions.sqrt_div_noiseBlockCount_le}
\leanok
\uses{BooleanAnalysis.ThresholdFunctions.noiseBlockCount,BooleanAnalysis.ThresholdFunctions.noiseBlockCount_spec}
\difficulty{5}
For every real $\delta$ with $0<\delta\leq1/2$, setting $m=\lfloor1/\delta\rfloor$ gives
\[
\frac{\sqrt m}{m}\leq\sqrt2\,\sqrt\delta.
\]
\end{lemma}

\begin{lemma}[A vanishing uniform sensitivity bound]
\lean{BooleanAnalysis.ThresholdFunctions.isUniformlyNoiseStable_of_bound}
\leanok
\uses{BooleanAnalysis.BooleanFunc,BooleanAnalysis.ThresholdFunctions.BooleanFunctionClass,BooleanAnalysis.ThresholdFunctions.IsBooleanValuedClass,BooleanAnalysis.ThresholdFunctions.IsUniformlyNoiseStable,BooleanAnalysis.ThresholdFunctions.noiseSensitivity,BooleanAnalysis.ThresholdFunctions.noiseSensitivity_le_half}
\difficulty{4}
Let $B=(B_n)_{n\in\mathbb N}$ be a family of sets of functions $\{0,1\}^{n}\to\{-1,1\}$ and let $\varepsilon:\mathbb R\to\mathbb R$ satisfy $\varepsilon(\delta)\geq0$ for $0\leq\delta\leq1/2$, $\lim_{\delta\to0^{+}}\varepsilon(\delta)=0$ and $\operatorname{NS}_{\delta}(f)\leq\varepsilon(\delta)$ for every $n$, $f\in B_n$ and $0<\delta\leq1/2$.
Then $B$ is uniformly noise-stable: there exists $\eta:\mathbb R\to\mathbb R$ with $0\leq\eta(\delta)\leq1$ on $[0,1/2]$, $\lim_{\delta\to0^{+}}\eta(\delta)=0$ and $\operatorname{NS}_{\delta}(f)\leq\eta(\delta)$ for every such $n,f,\delta$, where $\operatorname{NS}_{\delta}(f)=\Pr[f(X)\neq f(Y)]$, $X$ is uniform on the cube and $Y$ flips each coordinate independently with probability $\delta$.
\end{lemma}

\begin{lemma}[Boolean values of linear thresholds]
\lean{BooleanAnalysis.ThresholdFunctions.linearThresholdClass_isBooleanValued}
\leanok
\uses{BooleanAnalysis.ThresholdFunctions.IsBooleanValuedClass,BooleanAnalysis.ThresholdFunctions.linearThresholdClass}
\difficulty{1}
For every $n\in\mathbb N$, every member $f$ of the linear threshold class on $\{0,1\}^{n}$ takes values in $\{-1,1\}$; the class consists of functions $f(x)=\operatorname{sgn}_{+}(c+\sum_{i=1}^{n}a_i(-1)^{x_i})$ with real coefficients, where $\operatorname{sgn}_{+}(t)=1$ for $t\geq0$ and $-1$ for $t<0$.
\end{lemma}

\begin{lemma}[Closure of linear threshold functions]
\lean{BooleanAnalysis.ThresholdFunctions.linearThresholdClass_closed}
\leanok
\uses{BooleanAnalysis.BoolCube,BooleanAnalysis.ThresholdFunctions.IsClosedUnderInputNegationAndIdentification,BooleanAnalysis.ThresholdFunctions.IsLinearThreshold.exists_affine,BooleanAnalysis.ThresholdFunctions.identifiedNegation,BooleanAnalysis.ThresholdFunctions.isLinearThreshold_affine,BooleanAnalysis.ThresholdFunctions.isPmOne_thresholdSign,BooleanAnalysis.ThresholdFunctions.linearThresholdClass,BooleanAnalysis.ThresholdFunctions.thresholdSign,BooleanAnalysis.boolToSign}
\difficulty{5}
For all $n,m\in\mathbb N$, real coefficients $c,a_1,\ldots,a_n$, $z\in\{0,1\}^{n}$ and maps $\pi:\{1,\ldots,n\}\to\{1,\ldots,m\}$, if $f(x)=\operatorname{sgn}_{+}(c+\sum_{i=1}^{n}a_i(-1)^{x_i})$, then $g(w)=f((z_i+w_{\pi(i)}\pmod2)_{i=1}^{n})$ is a $\{-1,1\}$-valued linear threshold function on $\{0,1\}^{m}$, meaning that $g(w)=\operatorname{sgn}_{+}(b+\sum_{j=1}^{m}u_j(-1)^{w_j})$ for some real $b,u_1,\ldots,u_m$.
Here $\operatorname{sgn}_{+}(t)=1$ for $t\geq0$ and $-1$ for $t<0$.
\end{lemma}

\begin{definition}[Odd-multiplicity image]
\lean{BooleanAnalysis.ThresholdFunctions.oddFiberImage}
\leanok
For $n,m\in\mathbb N$, a map $\pi:\{1,\ldots,n\}\to\{1,\ldots,m\}$ and a subset $S\subseteq\{1,\ldots,n\}$, the odd-multiplicity image is the set
\[
O_{\pi}(S)=\{j\in\{1,\ldots,m\}: |\{i\in S:\pi(i)=j\}|\text{ is odd}\}.
\]
\end{definition}

\begin{lemma}[Cardinality of the odd-multiplicity image]
\lean{BooleanAnalysis.ThresholdFunctions.oddFiberImage_card_le}
\leanok
\uses{BooleanAnalysis.ThresholdFunctions.oddFiberImage,BooleanAnalysis.ThresholdFunctions.polynomialThresholdClass_mono}
\difficulty{4}
For all $n,m\in\mathbb N$, maps $\pi:\{1,\ldots,n\}\to\{1,\ldots,m\}$ and subsets $S\subseteq\{1,\ldots,n\}$, the set $O_{\pi}(S)=\{j\in\{1,\ldots,m\}:|\{i\in S:\pi(i)=j\}|\text{ is odd}\}$ satisfies $|O_{\pi}(S)|\leq|S|$.
\end{lemma}

\begin{lemma}[Parity of a sign power]
\lean{BooleanAnalysis.ThresholdFunctions.boolToSign_pow}
\leanok
\uses{BooleanAnalysis.boolToSign}
\difficulty{5}
For every bit $b\in\{0,1\}$ and natural number $c$, setting $s(b)=(-1)^b$ gives
\[
s(b)^c=\begin{cases}s(b)&\text{if }c\text{ is odd},\\1&\text{if }c\text{ is even}.\end{cases}
\]
\end{lemma}

\begin{lemma}[Sign products after coordinate identification]
\lean{BooleanAnalysis.ThresholdFunctions.prod_boolToSign_comp}
\leanok
\uses{BooleanAnalysis.BoolCube,BooleanAnalysis.ThresholdFunctions.boolToSign_pow,BooleanAnalysis.ThresholdFunctions.oddFiberImage,BooleanAnalysis.boolToSign}
\difficulty{4}
For $n,m\in\mathbb N$, a map $\pi:\{1,\ldots,n\}\to\{1,\ldots,m\}$, $S\subseteq\{1,\ldots,n\}$ and $w\in\{0,1\}^{m}$, set $O_{\pi}(S)=\{j\in\{1,\ldots,m\}:|\{i\in S:\pi(i)=j\}|\text{ is odd}\}$; then
\[
\prod_{i\in S}(-1)^{w_{\pi(i)}}=\prod_{j\in O_{\pi}(S)}(-1)^{w_j}.
\]
\end{lemma}

\begin{lemma}[Walsh characters under negated identification]
\lean{BooleanAnalysis.ThresholdFunctions.chiS_identifiedNegation}
\leanok
\uses{BooleanAnalysis.BoolCube,BooleanAnalysis.ThresholdFunctions.oddFiberImage,BooleanAnalysis.ThresholdFunctions.prod_boolToSign_comp,BooleanAnalysis.boolToSign,BooleanAnalysis.chiS}
\difficulty{4}
For $n,m\in\mathbb N$, $z\in\{0,1\}^{n}$, $w\in\{0,1\}^{m}$, a map $\pi:\{1,\ldots,n\}\to\{1,\ldots,m\}$ and $S\subseteq\{1,\ldots,n\}$, define $\chi_A(x)=\prod_{i\in A}(-1)^{x_i}$ for subsets of the coordinate set of $x$, and set $O_{\pi}(S)=\{j\in\{1,\ldots,m\}:|\{i\in S:\pi(i)=j\}|\text{ is odd}\}$; then
\[
\chi_S\bigl((z_i+w_{\pi(i)}\pmod2)_{i=1}^{n}\bigr)=\chi_S(z)\chi_{O_{\pi}(S)}(w).
\]
\end{lemma}

\begin{lemma}[Closure of polynomial threshold functions]
\lean{BooleanAnalysis.ThresholdFunctions.polynomialThresholdClass_closed}
\leanok
\uses{BooleanAnalysis.BoolCube,BooleanAnalysis.ThresholdFunctions.IsClosedUnderInputNegationAndIdentification,BooleanAnalysis.ThresholdFunctions.MultilinearPolynomial,BooleanAnalysis.ThresholdFunctions.MultilinearPolynomial.HasDegreeAtMost,BooleanAnalysis.ThresholdFunctions.MultilinearPolynomial.eval,BooleanAnalysis.ThresholdFunctions.MultilinearPolynomial.threshold,BooleanAnalysis.ThresholdFunctions.chiS_identifiedNegation,BooleanAnalysis.ThresholdFunctions.identifiedNegation,BooleanAnalysis.ThresholdFunctions.isPmOne_thresholdSign,BooleanAnalysis.ThresholdFunctions.oddFiberImage,BooleanAnalysis.ThresholdFunctions.oddFiberImage_card_le,BooleanAnalysis.ThresholdFunctions.polynomialThresholdClass,BooleanAnalysis.ThresholdFunctions.thresholdSign,BooleanAnalysis.chiS}
\difficulty{6}
For every $k,n,m\in\mathbb N$, a function $f:\{0,1\}^{n}\to\{-1,1\}$ of the form $f(x)=\operatorname{sgn}_{+}(p((-1)^{x_1},\ldots,(-1)^{x_n}))$ for a real multilinear polynomial $p$ of degree at most $k$, every $z\in\{0,1\}^{n}$ and every map $\pi:\{1,\ldots,n\}\to\{1,\ldots,m\}$, the function $g(w)=f((z_i+w_{\pi(i)}\pmod2)_{i=1}^{n})$ is also $\{-1,1\}$-valued and has a representation $g(w)=\operatorname{sgn}_{+}(q((-1)^{w_1},\ldots,(-1)^{w_m}))$ with a real multilinear polynomial $q$ of degree at most $k$.
Here $\operatorname{sgn}_{+}(t)=1$ for $t\geq0$ and $-1$ for $t<0$.
\end{lemma}

\begin{lemma}[Monotonicity in the allowed polynomial degree]
\lean{BooleanAnalysis.ThresholdFunctions.polynomialThresholdClass_mono}
\leanok
\uses{BooleanAnalysis.ThresholdFunctions.polynomialThresholdClass}
\difficulty{4}
For every $n,k,k'\in\mathbb N$ with $k\leq k'$, every $\{-1,1\}$-valued function on $\{0,1\}^{n}$ represented by $\operatorname{sgn}_{+}(p((-1)^{x_1},\ldots,(-1)^{x_n}))$ for a real multilinear polynomial $p$ of degree at most $k$ also admits such a representation of degree at most $k'$, where $\operatorname{sgn}_{+}(t)=1$ for $t\geq0$ and $-1$ for $t<0$.
\end{lemma}

\begin{lemma}[Exclusive-or form of negated identification]
\lean{BooleanAnalysis.ThresholdFunctions.identifiedNegation_apply}
\leanok
\uses{BooleanAnalysis.BoolCube,BooleanAnalysis.BooleanFunc,BooleanAnalysis.ThresholdFunctions.identifiedNegation,BooleanAnalysis.ThresholdFunctions.xorVec}
\difficulty{4}
For every $n,m\in\mathbb N$, $f:\{0,1\}^{n}\to\mathbb R$, $z\in\{0,1\}^{n}$, $w\in\{0,1\}^{m}$ and map $\pi:\{1,\ldots,n\}\to\{1,\ldots,m\}$, define $g(v)=f((u_i(v))_{i=1}^{n})$, where $u_i(v)=1-v_{\pi(i)}$ when $z_i=1$ and $u_i(v)=v_{\pi(i)}$ when $z_i=0$; then
\[
g(w)=f\bigl(z\oplus(w_{\pi(i)})_{i=1}^{n}\bigr),
\]
where $\oplus$ is coordinatewise addition modulo $2$.
\end{lemma}

\begin{lemma}[Flipping an identified variable]
\lean{BooleanAnalysis.ThresholdFunctions.identifiedNegation_flip}
\leanok
\uses{BooleanAnalysis.BoolCube,BooleanAnalysis.BooleanFunc,BooleanAnalysis.ThresholdFunctions.fiber,BooleanAnalysis.ThresholdFunctions.flipSet,BooleanAnalysis.ThresholdFunctions.identifiedNegation,BooleanAnalysis.ThresholdFunctions.identifiedNegation_apply,BooleanAnalysis.ThresholdFunctions.mem_fiber,BooleanAnalysis.ThresholdFunctions.xorVec,BooleanAnalysis.flipBit}
\difficulty{5}
For $n,m\in\mathbb N$, $f:\{0,1\}^{n}\to\mathbb R$, $z\in\{0,1\}^{n}$, a map $\pi:\{1,\ldots,n\}\to\{1,\ldots,m\}$, $w\in\{0,1\}^{m}$ and $j\in\{1,\ldots,m\}$, define $g(v)=f(z\oplus(v_{\pi(i)})_{i=1}^{n})$, where $\oplus$ is coordinatewise addition modulo $2$, and write $x^{A}$ for the vector obtained by flipping precisely the coordinates in $A$; then
\[
g(w^{\{j\}})=f\bigl((z\oplus(w_{\pi(i)})_{i=1}^{n})^{\pi^{-1}(\{j\})}\bigr).
\]
\end{lemma}

\begin{lemma}[Average influence of an identified variable]
\lean{BooleanAnalysis.ThresholdFunctions.expect_influence_identifiedNegation}
\leanok
\uses{BooleanAnalysis.BoolCube,BooleanAnalysis.BooleanFunc,BooleanAnalysis.ThresholdFunctions.fiber,BooleanAnalysis.ThresholdFunctions.flipDisagreement,BooleanAnalysis.ThresholdFunctions.flipSet,BooleanAnalysis.ThresholdFunctions.identifiedNegation,BooleanAnalysis.ThresholdFunctions.identifiedNegation_apply,BooleanAnalysis.ThresholdFunctions.identifiedNegation_flip,BooleanAnalysis.ThresholdFunctions.sum_comp_xorVec,BooleanAnalysis.ThresholdFunctions.xorVec,BooleanAnalysis.expect,BooleanAnalysis.expect_eq_fintypeExpect,BooleanAnalysis.flipBit,BooleanAnalysis.influence,BooleanAnalysis.uniformWeight}
\difficulty{5}
For $n,m\in\mathbb N$, $f:\{0,1\}^{n}\to\mathbb R$, a map $\pi:\{1,\ldots,n\}\to\{1,\ldots,m\}$ and $j\in\{1,\ldots,m\}$, define $g_z(w)=f(z\oplus(w_{\pi(i)})_{i=1}^{n})$, where $\oplus$ is coordinatewise addition modulo $2$, and write $x^{A}$ for the vector obtained by flipping the coordinates in $A$; then
\[
\mathbb E_{z}\!\left[\operatorname{Inf}_j(g_z)\right]
=\mathbb E_x\!\left[\frac{(f(x)-f(x^{\pi^{-1}(\{j\})}))^2}{4}\right],
\qquad
\operatorname{Inf}_j(h)=\mathbb E_w\!\left[\frac{(h(w)-h(w^{\{j\}}))^2}{4}\right].
\]
Each expectation is uniform on the full Boolean cube of its variable's dimension.
\end{lemma}

\begin{lemma}[Sum of block-flip disagreements]
\lean{BooleanAnalysis.ThresholdFunctions.sum_flipDisagreement_fiber}
\leanok
\uses{BooleanAnalysis.BooleanFunc,BooleanAnalysis.ThresholdFunctions.fiber,BooleanAnalysis.ThresholdFunctions.flipDisagreement,BooleanAnalysis.ThresholdFunctions.flipDisagreement_eq,BooleanAnalysis.ThresholdFunctions.noiseSensitivity,BooleanAnalysis.ThresholdFunctions.noiseStability,BooleanAnalysis.ThresholdFunctions.sum_pi_sign,BooleanAnalysis.fourierCoeff,BooleanAnalysis.isPmOne,BooleanAnalysis.stability_formula}
\difficulty{5}
For every $n\in\mathbb N$, positive integer $m$, $f:\{0,1\}^{n}\to\{-1,1\}$ and $j\in\{1,\ldots,m\}$,
\[
\sum_{\pi:\{1,\ldots,n\}\to\{1,\ldots,m\}}
\Pr_X\!\left[f(X)\neq f(X^{\pi^{-1}(\{j\})})\right]
=m^n\operatorname{NS}_{1/m}(f).
\]
Here $X$ is uniform on $\{0,1\}^{n}$, $X^{A}$ flips precisely the coordinates of $A$, and $\operatorname{NS}_{1/m}(f)$ is $\Pr[f(X)\neq f(Y)]$ when each coordinate of $X$ is independently flipped with probability $1/m$ to obtain $Y$.
\end{lemma}

\begin{lemma}[Monotonicity of totalized noise sensitivity]
\lean{BooleanAnalysis.ThresholdFunctions.noiseSensitivity_mono}
\leanok
\uses{BooleanAnalysis.BooleanFunc,BooleanAnalysis.ThresholdFunctions.noiseSensitivity,BooleanAnalysis.ThresholdFunctions.noiseStability,BooleanAnalysis.stability_formula}
\difficulty{4}
For every $n\in\mathbb N$, $f:\{0,1\}^{n}\to\mathbb R$ and real $\delta\leq\delta'\leq1/2$, with no lower bound on either rate, one has $\operatorname{NS}_{\delta}(f)\leq\operatorname{NS}_{\delta'}(f)$, where the total real-valued extension is
\[
\operatorname{NS}_{t}(f)=\frac{1-\sum_{S\subseteq\{1,\ldots,n\}}(1-2t)^{|S|}\widehat f(S)^2}{2},
\qquad
\widehat f(S)=2^{-n}\sum_{x\in\{0,1\}^{n}}f(x)\prod_{i\in S}(-1)^{x_i}.
\]
\end{lemma}

\begin{theorem}[Influence-to-noise bound for closed classes]
\lean{BooleanAnalysis.ThresholdFunctions.influence_bound_implies_noise_bound}
\leanok
\uses{BooleanAnalysis.BoolCube,BooleanAnalysis.BooleanFunc,BooleanAnalysis.ThresholdFunctions.BooleanFunctionClass,BooleanAnalysis.ThresholdFunctions.IsBooleanValuedClass,BooleanAnalysis.ThresholdFunctions.IsClosedUnderInputNegationAndIdentification,BooleanAnalysis.ThresholdFunctions.expect_influence_identifiedNegation,BooleanAnalysis.ThresholdFunctions.identifiedNegation,BooleanAnalysis.ThresholdFunctions.noiseBlockCount,BooleanAnalysis.ThresholdFunctions.noiseBlockCount_spec,BooleanAnalysis.ThresholdFunctions.noiseSensitivity,BooleanAnalysis.ThresholdFunctions.noiseSensitivity_mono,BooleanAnalysis.ThresholdFunctions.sum_flipDisagreement_fiber,BooleanAnalysis.expect,BooleanAnalysis.expect_eq_fintypeExpect,BooleanAnalysis.influence,BooleanAnalysis.totalInfluence}
\difficulty{6}
Let $B=(B_n)_{n\in\mathbb N}$ be a family of sets of functions $\{0,1\}^{n}\to\{-1,1\}$, closed under all substitutions $w\mapsto(z_i+w_{\pi(i)}\pmod2)_{i=1}^{n}$ for every $n,m\in\mathbb N$, $z\in\{0,1\}^{n}$ and map $\pi:\{1,\ldots,n\}\to\{1,\ldots,m\}$, and let $A:\mathbb N\to\mathbb R$ satisfy
\[
\operatorname{Inf}(f):=\sum_{i=1}^{n}\Pr_X[f(X)\neq f(X^{\{i\}})]\leq A(n)
\quad\text{for every }n\in\mathbb N,\ f\in B_n,
\]
where $X$ is uniform on $\{0,1\}^{n}$ and $X^{\{i\}}$ flips coordinate $i$.
Then for every $0<\delta\leq1/2$, $n\in\mathbb N$ and $f\in B_n$, setting $m=\lfloor1/\delta\rfloor$ gives
\[
\operatorname{NS}_{\delta}(f)\leq\frac{A(m)}{m},
\]
where $\operatorname{NS}_{\delta}(f)=\Pr[f(X)\neq f(Y)]$ and $Y$ independently flips each coordinate of $X$ with probability $\delta$.
\end{theorem}

\begin{theorem}[Peres's noise-sensitivity bound]
\lean{BooleanAnalysis.ThresholdFunctions.peres_theorem}
\leanok
\uses{BooleanAnalysis.BooleanFunc,BooleanAnalysis.ThresholdFunctions.influence_bound_implies_noise_bound,BooleanAnalysis.ThresholdFunctions.linearThresholdClass,BooleanAnalysis.ThresholdFunctions.linearThresholdClass_closed,BooleanAnalysis.ThresholdFunctions.linearThresholdClass_isBooleanValued,BooleanAnalysis.ThresholdFunctions.ltf_totalInfluence_le_sqrt,BooleanAnalysis.ThresholdFunctions.noiseSensitivity,BooleanAnalysis.ThresholdFunctions.sqrt_div_noiseBlockCount_le}
\difficulty{3}
There exists a universal real constant $C\geq0$ such that, for every $n\in\mathbb N$, every linear threshold function $f(x)=\operatorname{sgn}_{+}(c+\sum_{i=1}^{n}a_i(-1)^{x_i})$ on $\{0,1\}^{n}$ with real coefficients and every $0<\delta\leq1/2$,
\[
\operatorname{NS}_{\delta}(f)\leq C\sqrt\delta.
\]
Here $\operatorname{sgn}_{+}(t)=1$ for $t\geq0$ and $-1$ for $t<0$, and $\operatorname{NS}_{\delta}(f)=\Pr[f(X)\neq f(Y)]$ with uniform $X$ and independent coordinate flips of probability $\delta$ producing $Y$.
\end{theorem}

\begin{theorem}[Uniform stability of linear thresholds]
\lean{BooleanAnalysis.ThresholdFunctions.linearThreshold_uniformlyNoiseStable}
\leanok
\uses{BooleanAnalysis.ThresholdFunctions.IsUniformlyNoiseStable,BooleanAnalysis.ThresholdFunctions.isUniformlyNoiseStable_of_bound,BooleanAnalysis.ThresholdFunctions.linearThresholdClass,BooleanAnalysis.ThresholdFunctions.linearThresholdClass_isBooleanValued,BooleanAnalysis.ThresholdFunctions.peres_theorem}
\difficulty{4}
The family of linear threshold functions on Boolean cubes is uniformly noise-stable: every function has the form $f(x)=\operatorname{sgn}_{+}(c+\sum_{i=1}^{n}a_i(-1)^{x_i})$ for real coefficients, is $\{-1,1\}$-valued, and there exists $\varepsilon:\mathbb R\to\mathbb R$ satisfying $0\leq\varepsilon(\delta)\leq1$ on $[0,1/2]$, $\lim_{\delta\to0^{+}}\varepsilon(\delta)=0$ and
\[
\Pr[f(X)\neq f(Y)]\leq\varepsilon(\delta)
\]
for every dimension $n\in\mathbb N$, every such $f$ and every $0<\delta\leq1/2$.
Here $\operatorname{sgn}_{+}(t)=1$ for $t\geq0$ and $-1$ for $t<0$, $X$ is uniform on $\{0,1\}^{n}$, and $Y$ independently flips its coordinates with probability $\delta$.
\end{theorem}

\begin{lemma}[Vanishing square-root polylogarithmic modulus]
\lean{BooleanAnalysis.ThresholdFunctions.tendsto_sqrt_mul_log_pow}
\leanok
\difficulty{5}
For every natural number $d$, with $\log$ denoting the natural logarithm,
\[
\lim_{\delta\to0^{+}}\sqrt\delta\,\bigl(\log(1/\delta)\bigr)^d=0.
\]
\end{lemma}

\begin{theorem}[Kane's influence bound for higher degrees]
\lean{BooleanAnalysis.ThresholdFunctions.kane_totalInfluence_bound_of_two_le}
\leanok
\uses{BooleanAnalysis.BooleanFunc,BooleanAnalysis.ThresholdFunctions.polynomialThresholdClass,BooleanAnalysis.totalInfluence}
There exists a universal real constant $C>0$ such that, for every $n,k\in\mathbb N$ with $k\geq2$ and every $f:\{0,1\}^{n}\to\{-1,1\}$ admitting a representation $f(x)=\operatorname{sgn}_{+}(p((-1)^{x_1},\ldots,(-1)^{x_n}))$ by a real multilinear polynomial $p$ of degree at most $k$,
\[
\operatorname{Inf}(f)\leq\sqrt n\,
\bigl(2^k\log(\max\{n,2\})\bigr)^{Ck\log(\max\{k,2\})}.
\]
Here $\operatorname{Inf}(f)=\sum_{i=1}^{n}\Pr_X[f(X)\neq f(X^{\{i\}})]$, $X$ is uniform, $X^{\{i\}}$ flips coordinate $i$, $\log$ is the natural logarithm, real powers have positive bases, and $\operatorname{sgn}_{+}(t)=1$ for $t\geq0$ and $-1$ for $t<0$.
\end{theorem}

\begin{lemma}[Degree-one polynomial thresholds are linear]
\lean{BooleanAnalysis.ThresholdFunctions.isLinearThreshold_of_polynomialThresholdClass}
\leanok
\uses{BooleanAnalysis.BooleanFunc,BooleanAnalysis.ThresholdFunctions.IsLinearThreshold,BooleanAnalysis.ThresholdFunctions.polynomialThresholdClass}
\difficulty{4}
For every $n,k\in\mathbb N$ with $k\leq1$, if $f:\{0,1\}^{n}\to\{-1,1\}$ has a representation $f(x)=\operatorname{sgn}_{+}(p((-1)^{x_1},\ldots,(-1)^{x_n}))$ by a real multilinear polynomial $p$ of degree at most $k$, then there exist real $c,a_1,\ldots,a_n$ such that
\[
f(x)=\operatorname{sgn}_{+}\left(c+\sum_{i=1}^{n}a_i(-1)^{x_i}\right)\quad\text{for all }x\in\{0,1\}^{n},
\]
where $\operatorname{sgn}_{+}(t)=1$ for $t\geq0$ and $-1$ for $t<0$.
\end{lemma}

\begin{theorem}[Kane's influence bound for all degrees]
\lean{BooleanAnalysis.ThresholdFunctions.kane_totalInfluence_le}
\leanok
\uses{BooleanAnalysis.BooleanFunc,BooleanAnalysis.ThresholdFunctions.isLinearThreshold_of_polynomialThresholdClass,BooleanAnalysis.ThresholdFunctions.kane_totalInfluence_bound_of_two_le,BooleanAnalysis.ThresholdFunctions.ltf_totalInfluence_le_sqrt,BooleanAnalysis.ThresholdFunctions.polynomialThresholdClass,BooleanAnalysis.totalInfluence}
\difficulty{5}
There exists a universal real constant $C>0$ such that, for every $n,k\in\mathbb N$, including $n=0$ and $k=0$, and every $f:\{0,1\}^{n}\to\{-1,1\}$ represented as $f(x)=\operatorname{sgn}_{+}(p((-1)^{x_1},\ldots,(-1)^{x_n}))$ by a real multilinear polynomial $p$ of degree at most $k$,
\[
\operatorname{Inf}(f)\leq\sqrt n\,
\bigl(2^k\log(\max\{n,2\})\bigr)^{Ck\log(\max\{k,2\})}.
\]
Here $\operatorname{Inf}(f)=\sum_{i=1}^{n}\Pr_X[f(X)\neq f(X^{\{i\}})]$, $X$ is uniform, $X^{\{i\}}$ flips coordinate $i$, $\log$ is the natural logarithm, real powers have positive bases, and $\operatorname{sgn}_{+}(t)=1$ for $t\geq0$ and $-1$ for $t<0$.
\end{theorem}

\begin{lemma}[Bounding a real logarithmic exponent]
\lean{BooleanAnalysis.ThresholdFunctions.rpow_le_log_two_rpow_mul_pow}
\leanok
\difficulty{5}
For real $L,E$ and $d\in\mathbb N$ with $L\geq\log2$ and $0\leq E\leq d$, where $\log$ is the natural logarithm, one has
\[
L^E\leq(\log2)^{-d}L^d,
\]
with real exponentiation on the positive bases $L$ and $\log2$.
\end{lemma}

\begin{lemma}[Block-count polylogarithmic estimate]
\lean{BooleanAnalysis.ThresholdFunctions.kane_noise_aux}
\leanok
\uses{BooleanAnalysis.ThresholdFunctions.noiseBlockCount,BooleanAnalysis.ThresholdFunctions.noiseBlockCount_spec,BooleanAnalysis.ThresholdFunctions.rpow_le_log_two_rpow_mul_pow,BooleanAnalysis.ThresholdFunctions.sqrt_div_noiseBlockCount_le}
\difficulty{5}
For $k,d\in\mathbb N$ and real $E,\delta$ with $0\leq E\leq d$ and $0<\delta\leq1/2$, setting $m=\lfloor1/\delta\rfloor$ gives
\[
\frac{\sqrt m\,\bigl(2^k\log(\max\{m,2\})\bigr)^E}{m}
\leq\sqrt2\,(2^k)^E(\log2)^{-d}\sqrt\delta\,\bigl(\log(1/\delta)\bigr)^d,
\]
where $\log$ is the natural logarithm and the powers with real exponents have positive bases.
\end{lemma}

\begin{theorem}[Kane's influence and noise-sensitivity bounds]
\lean{BooleanAnalysis.ThresholdFunctions.kane_totalInfluence_bound}
\leanok
\uses{BooleanAnalysis.BooleanFunc,BooleanAnalysis.ThresholdFunctions.influence_bound_implies_noise_bound,BooleanAnalysis.ThresholdFunctions.kane_noise_aux,BooleanAnalysis.ThresholdFunctions.kane_totalInfluence_le,BooleanAnalysis.ThresholdFunctions.noiseSensitivity,BooleanAnalysis.ThresholdFunctions.polynomialThresholdClass,BooleanAnalysis.ThresholdFunctions.polynomialThresholdClass_closed,BooleanAnalysis.totalInfluence}
\difficulty{5}
For $n,k\in\mathbb N$, let a degree-at-most-$k$ polynomial threshold function mean $f:\{0,1\}^{n}\to\{-1,1\}$ of the form $f(x)=\operatorname{sgn}_{+}(p((-1)^{x_1},\ldots,(-1)^{x_n}))$ for a real multilinear polynomial $p$ of degree at most $k$, with $\operatorname{sgn}_{+}(t)=1$ for $t\geq0$ and $-1$ for $t<0$; there exists a universal $C_0>0$ such that every such function satisfies
\[
\operatorname{Inf}(f)\leq\sqrt n\,
\bigl(2^k\log(\max\{n,2\})\bigr)^{C_0k\log(\max\{k,2\})}.
\]
In addition, for every fixed $k\in\mathbb N$ with $k>0$, there exist a real $C_k\geq0$ and a natural number $d_k$ such that every dimension $n\in\mathbb N$, every degree-at-most-$k$ polynomial threshold function $f$ and every $0<\delta\leq1/2$ satisfy
\[
\operatorname{NS}_{\delta}(f)\leq C_k\sqrt\delta\,\bigl(\log(1/\delta)\bigr)^{d_k}.
\]
Here $\operatorname{Inf}(f)=\sum_{i=1}^{n}\Pr_X[f(X)\neq f(X^{\{i\}})]$, $X$ is uniform, $X^{\{i\}}$ flips coordinate $i$, $\operatorname{NS}_{\delta}(f)=\Pr[f(X)\neq f(Y)]$ when $Y$ independently flips each coordinate of $X$ with probability $\delta$, and $\log$ is the natural logarithm.
\end{theorem}

\begin{theorem}[Uniform stability of polynomial thresholds]
\lean{BooleanAnalysis.ThresholdFunctions.polynomialThreshold_uniformlyNoiseStable}
\leanok
\uses{BooleanAnalysis.ThresholdFunctions.IsUniformlyNoiseStable,BooleanAnalysis.ThresholdFunctions.isUniformlyNoiseStable_of_bound,BooleanAnalysis.ThresholdFunctions.kane_totalInfluence_bound,BooleanAnalysis.ThresholdFunctions.polynomialThresholdClass,BooleanAnalysis.ThresholdFunctions.polynomialThresholdClass_mono,BooleanAnalysis.ThresholdFunctions.tendsto_sqrt_mul_log_pow}
\difficulty{5}
For every fixed $k\in\mathbb N$, including $k=0$, the degree-at-most-$k$ polynomial threshold functions form a uniformly noise-stable family: there exists $\varepsilon:\mathbb R\to\mathbb R$ with $0\leq\varepsilon(\delta)\leq1$ for $0\leq\delta\leq1/2$ and $\lim_{\delta\to0^{+}}\varepsilon(\delta)=0$ such that
\[
\Pr[f(X)\neq f(Y)]\leq\varepsilon(\delta)
\]
for every $n\in\mathbb N$, $0<\delta\leq1/2$ and $f:\{0,1\}^{n}\to\{-1,1\}$ admitting a representation $f(x)=\operatorname{sgn}_{+}(p((-1)^{x_1},\ldots,(-1)^{x_n}))$ with a real multilinear polynomial $p$ of degree at most $k$.
Here $\operatorname{sgn}_{+}(t)=1$ for $t\geq0$ and $-1$ for $t<0$, $X$ is uniform on $\{0,1\}^{n}$ and $Y$ independently flips each coordinate of $X$ with probability $\delta$.
\end{theorem}
